\chapter{A Map of the Series}\label{ch:map}

\emph{Order and Chance} is a road from the mathematics of a first university course to the research frontier of
three subjects that turn out to be one: number theory, algebra and probability. The destination is probabilistic
number theory, the study of the integers as if they were random, which needs all three. This chapter describes the
shape of the road.

\section{The parts}

\textbf{Part I, Foundations} (Tomes 0 to 5), rebuilds the toolkit and then the three subjects. Tome~0 is a
problem-driven review of undergraduate mathematics: sums, complex numbers, polynomials, inequalities, series,
integrals, linear algebra and counting. Tome~1 is the arithmetic of the integers, from divisibility to quadratic
reciprocity, continued fractions and the first steps of algebraic number theory. Tome~2 is algebra: groups, rings,
fields, Galois theory and characters. Tome~3 is probability, from counting to martingales. Tome~4 is the art of
calculation: series, integrals, special functions and asymptotics. Tome~5 is where the threads first meet: analytic
and probabilistic number theory, from Dirichlet's theorem to the prime number theorem and the Erd\H os--Kac theorem.

\textbf{Part II, the graduate toolbox} (Tomes 6 to 9), supplies the machinery that research papers take for granted:
topology, measure and integration (Tome~6); functional, harmonic and complex analysis (Tome~7); manifolds, algebraic
topology and Lie groups (Tome~8); categories, commutative and homological algebra (Tome~9).

\textbf{Part III, toward research} (Tomes 10 to 15), climbs to the frontier: limit theorems, large deviations and
ergodic theory (Tome~10); Brownian motion, stochastic calculus and random matrices (Tome~11); algebraic number theory,
elliptic curves and modular forms (Tome~12); analytic and additive number theory (Tome~13); probabilistic number theory
and arithmetic statistics, the destination (Tome~14); and a workshop on reading papers, experimenting, conjecturing and
writing (Tome~15).

Beside the main road run three shorter ones. The \textbf{computing track} (Tomes C1 to C5) treats the mathematics of
computer science in the same spirit: discrete structures, algorithms, computation and complexity, information and
codes, optimisation and learning. The \textbf{companion volumes} (S1 to S4) serve the whole series: linear algebra in
depth, geometry, a book of open problems, and this map. The \textbf{problem books} (P1 to P3), one per part, mix the
tomes of a part in harder problems where nobody says which chapter applies.

\section{The tomes at a glance}

The table lists every tome of the plan, with the number of its chapters, of its named results, and of its exercises,
as they stand in this edition.

% Generated by generate.py: do not edit.
\begin{center}\small\adjustbox{max width=\linewidth,max totalheight=0.85\textheight}{%
\begin{tabular}{@{}lp{0.52\linewidth}rrr@{}}
Tome&Title&ch.&named&ex.\\\hline
\multicolumn{5}{@{}l}{\textit{Part I. Foundations}}\\
0&Back in Shape&8&69&130\\
1&The Integers&12&44&97\\
2&Structures&12&56&96\\
3&Chance&12&105&94\\
4&The Art of Calculation&12&94&72\\
5&Where It All Meets&12&59&83\\
\multicolumn{5}{@{}l}{\textit{Part II. The graduate toolbox}}\\
6&Spaces&11&70&83\\
7&Analysis in Depth&10&91&69\\
8&Shapes&10&78&62\\
9&Structures II&11&69&66\\
\multicolumn{5}{@{}l}{\textit{Part III. Toward research}}\\
10&Chance II&12&114&80\\
11&Random Motion&12&110&95\\
12&Arithmetic&14&103&110\\
13&Primes in Depth&14&101&80\\
14&Random Arithmetic&13&52&77\\
15&The Research Workshop&12&40&75\\
\multicolumn{5}{@{}l}{\textit{The Computing track}}\\
C1&Discrete Structures&12&114&136\\
C2&Algorithms&12&108&120\\
C3&Computation and Complexity&12&103&122\\
C4&Information, Codes and Secrets&4&28&43\\
C5&Data, Optimisation and Learning&2&29&24\\
\multicolumn{5}{@{}l}{\textit{Companion volumes}}\\
S1&Linear Algebra in Depth&12&107&147\\
S2&Geometry&12&72&143\\
S3&A Book of Conjectures&12&34&117\\
S4&The Companion (this book)&&&\\
\multicolumn{5}{@{}l}{\textit{Problem books}}\\
P1&Crossroads&14&0&300\\
P2&Machinery&2&0&44\\
P3&Frontiers&1&0&22\\
\end{tabular}}
\end{center}


\begin{think}
Before turning to the figure of the next section, guess: which tome of Part~I is cited most often by the rest of the
series? Which tome of Part~III leans most heavily on Part~II?
\end{think}

\section{Which tome leans on which}

Every tome states what it assumes, and every time it uses a result proved elsewhere it says where: ``Tome~5,
Theorem~4.1'', ``Tome~2, Chapter~10''. Counting these references draws the dependency map of the series
(Figure~\ref{fig:citations}). Read a row to see what a tome uses, and a column to see who uses it.

\begin{figure}[ht]
\centering
\begin{minipage}{0.85\linewidth}\centering\fig{citations}\end{minipage}
\caption{References between tomes: row $A$, column $B$ counts the places where Tome~$A$ cites a numbered result or a
chapter of Tome~$B$.}
\label{fig:citations}
\end{figure}

A few patterns stand out. The problem book P1 cites all of Part~I, which is its purpose. Tome~5 leans on Tome~4: the
analytic theory of numbers is built from the series, integrals and asymptotics of the art of calculation. Tome~7 is
built on Tome~6, and Tome~12 on Tomes~1 and~2. The probability tomes form their own chain, from Tome~3 to Tome~10 and
Tome~11. And Tome~S3, the book of conjectures, draws on Tome~5 and Tome~12 more than on anything else, because the open
problems of number theory live there.

\section{Reading paths}

The tomes are numbered in the order of the main road, and each one names its prerequisites in its preface. Three
other orders are natural.
\begin{itemize}
\item \emph{The shortest road to probabilistic number theory}: Tomes 0, 1, 3, 4, 5, then 6 and 7 for measure and
analysis, 10 for limit theorems, and 13 for the primes. Tomes 2, 9 and 12 add the algebraic side.
\item \emph{The algebraic road}: Tomes 0, 1, 2, then 9 and 12, with S1 and S2 alongside.
\item \emph{The computing road}: Tome~0, then C1 to C5, borrowing Tome~3 for probability and S1 for linear algebra.
\end{itemize}
Whatever the road, the companion volumes can be opened at any time: S1 after Tome~0, S2 after S1 and Tome~2, S3 after
Part~I, and this book whenever you need to find something.
